% Requiere \usepackage{booktabs}
\begin{table}[htbp]
  \centering\small
  \begin{tabular}{llrrrrrrrrrllrrrrrrl}
    \toprule
    variante & dtype & M & N & K & BLOCK\_SIZE\_M & BLOCK\_SIZE\_N & BLOCK\_SIZE\_K & GROUP\_SIZE\_M & num\_warps & num\_stages & carga & usa\_tma & autotune\_s & ms & ms\_p20 & ms\_p80 & tflops & gbs & mma \\
    \midrule
    baseline & fp16 & 8192 & 8192 & 8192 & 128 & 256 & 64 & 8 & 8 & 3 & cp.async & False & 13.9 & 12.1083 & 11.9548 & 12.554 & 90.81 & 33.3 & mma.sync.aligned.m16n8k16.row.col.f32.f16.f16.f32 \\
    blockptr & fp16 & 8192 & 8192 & 8192 & 128 & 256 & 64 & 8 & 8 & 3 & cp.async & False & 78.1 & 11.8462 & 11.8164 & 12.3852 & 92.82 & 34 & mma.sync.aligned.m16n8k16.row.col.f32.f16.f16.f32 \\
    tma & fp16 & 8192 & 8192 & 8192 & 128 & 128 & 32 & 8 & 4 & 4 & TMA & True & 97.7 & 11.8599 & 11.8177 & 12.2311 & 92.71 & 34 & mma.sync.aligned.m16n8k16.row.col.f32.f16.f16.f32 \\
    \bottomrule
  \end{tabular}
  \caption{experimento\_tma: NVIDIA GB10 (sm\_121), torch 2.9.0a0+145a3a7bda.nv25.10, triton 3.8.0, CUDA 13.0; job 20398, 2026-09-29T20:56:26. Mediana de triton.testing.do\_bench.}
  \label{tab:experimento-tma}
\end{table}
